\scanpage{3}
Si $f$ est un endomorphisme d'un vectoriel $E$ de dim finie $n$ sur le corps $k$, on pose
\[
\begin{aligned}
P_f(t)=\det(1-ft)&=\prod_i(1-\lambda_i t)\\
&=\sum_i(-1)^i\bigl(\operatorname{Tr}\lambda^i(f)\bigr)t^i
\end{aligned}
\]
les $\lambda_i$ étant les val propres dans une \unsure{extension|clôture} alg. close de $k$.

Donc les racines de $P_f(t)$ sont les inverses des valeurs propres non nulles de $f$.

$P_f(t)$ pour $(E,f)$ variable est ``additif'' en $E$ (\unsure{suite exacte|suites exactes} en $E$ \unsure{transformée|transformé} en produit des polynômes). Il en est donc de $=$ de
\[
\begin{aligned}
Z_f(t)=\frac{1}{P_f(t)}
 &=\prod_i\frac{1}{1-\lambda_i t}
 =\prod_i\sum_{\unsure{\nu|\kappa}}\lambda_i^\nu t^\nu\\
 &=\sum_\nu t^\nu\sum_{\nu_1+\cdots+\nu_n=\nu}\lambda_1^{\unsure{\nu_1|\nu_n}}\cdots\lambda_n^{\nu_n}
 =\sum_\nu\operatorname{Tr}\bigl(\operatorname{Sym}^\nu(f)\bigr)t^\nu
\end{aligned}
\]

D'autre part, $P_f(t)$ transforme le produit tensoriel en $(E,f)$ en le ``produit'' $\ast$ des séries formelles commençant par 1, où on a posé
\[
(1-at)\ast(1-bt)=1-abt
\]
[et distributivité de $\ast$ p.r. à\ $\cdot$, etc\ldots]
\margin{\diagram{\deleted{$\begin{array}{ccc}
p-1 & \underset{\unsure{\leftrightarrow|{-}}}{p} & p+1\\
\underset{\rightarrow}{-2} & -1 & \underset{\unsure{\leftarrow|\leftrightarrow}}{0}\\
& p,\ p+1 & \unsure{p+2|\rho+2}\\
p-1 & 0 & p+1
\end{array}$}}}
% archive sheet number: 2

